\chapter{Failure Taxonomy}
\label{chap:failure_taxonomy}

\section{Failure Mode Taxonomy}
\label{sec:ft_taxonomy}

Chapter~\ref{chap:results} reported \emph{how often} the agent fails.
This chapter is about \emph{how} it fails, and it is the part of the study we
regard as the original contribution.
Pass rates on formal benchmarks are widely published; what is not published, as
far as we have been able to establish, is a mechanically checkable taxonomy of
the ways a language-model proof agent goes wrong, applied to a complete run and
reported with its own error bars and its own failures of fit.
The taxonomy below is offered in that spirit: not as a finished instrument, but
as one specific enough to be wrong in public.

Five modes are defined:

\begin{itemize}
  \item \texttt{NO\_PROGRESS\_LOOP} --- Lean accepts the tactic but the goal
        does not change, repeatedly.
  \item \texttt{ERROR\_LOOP} --- Lean returns the identical message on three
        consecutive steps that produced a message at all.
  \item \texttt{HALLUCINATED\_LEMMA} --- the tactic names a capitalised symbol
        Lean does not know.
  \item \texttt{SYNTAX\_ERROR\_LOOP} --- three or more steps emit text Lean
        cannot parse.
  \item \texttt{BUDGET\_EXHAUSTION\_MIXED} --- none of the above.
\end{itemize}

The taxonomy is \emph{mutually exclusive} by construction: classification is
first-match-wins in the order listed, so a proof that would satisfy two tests
is assigned the earlier one. It is \emph{exhaustive} by construction too, but
only in the cheap way --- the fifth mode is a residual that accepts anything the
first four reject. Exhaustiveness bought like that is not evidence the taxonomy
carves the space well, and \S\ref{sec:ft_negative} reports that it does not.

Ordering is a modelling choice with consequences. \texttt{NO\_PROGRESS\_LOOP}
is tested first because it is the most specific claim; \texttt{ERROR\_LOOP}
precedes \texttt{HALLUCINATED\_LEMMA} because a proof that both loops on an
error and invents a lemma is better described by the loop. A different order
would produce different counts from identical logs. The counts in this chapter
are therefore properties of \texttt{classify\_failures.py} as much as of the
agent, which is the honest way to read any taxonomy of this kind.

\section{Detection Rules}
\label{sec:ft_rules}

Each rule below is the actual code from \texttt{classify\_failures.py}, not a
paraphrase. A proof is represented by \texttt{proof.steps}, the step events in
order; \texttt{\_first\_message(s)} returns the first entry of that step's
\texttt{lean\_messages} or \texttt{None}; \texttt{\_joined(s)} returns all of
them lowercased and concatenated.

\subsection{NO\_PROGRESS\_LOOP}

\begin{verbatim}
    if len(steps) >= 3:
        tail = steps[-3:]
        if all(s.get("outcome") == "no_progress" for s in tail):
            return "NO_PROGRESS_LOOP", tail
\end{verbatim}

\noindent
\textbf{Catches:} an agent whose last three tactics were all accepted by Lean
without changing the goal --- \texttt{skip}-like tactics, a \texttt{simp} with
nothing to rewrite, a \texttt{try} whose body failed silently.

\noindent
\textbf{Misses:} any such loop that is interrupted once in its final three
steps, and --- decisively for this study --- any agent that never receives the
\texttt{no\_progress} outcome in the first place. See \S\ref{sec:ft_zero}.

\subsection{ERROR\_LOOP}

\begin{verbatim}
    with_msgs = [s for s in steps if _first_message(s)]
    if len(with_msgs) >= 3:
        tail = with_msgs[-3:]
        first = _first_message(tail[0])
        if all(_first_message(s) == first for s in tail[1:]):
            return "ERROR_LOOP", tail
\end{verbatim}

\noindent
\textbf{Catches:} exact repetition of Lean's first message across the last
three message-producing steps. Note that the filter is applied before the tail
is taken, so silent steps in between do not break the loop.

\noindent
\textbf{Misses:} semantically identical errors with different text --- two
different unknown identifiers, or the same type mismatch reported against
different metavariable numbers, will not compare equal. Conversely it
\emph{over}-catches: the rule says nothing about where the message came from,
so three identical provider errors satisfy it exactly. That is the subject of
\S\ref{sec:ft_provider}, and it turns out to account for 22 of 28 raw hits.

\subsection{HALLUCINATED\_LEMMA}

\begin{verbatim}
    hallucinated = [s for s in steps
                    if "unknown identifier" in _joined(s)
                    and names_a_hallucinated_symbol(s.get("tactic"))]
    if hallucinated:
        return "HALLUCINATED_LEMMA", hallucinated[:3]
\end{verbatim}

\noindent where the helper is

\begin{verbatim}
    STDLIB_NAMES = frozenset({
        "Nat", "Int", "Real", "List", "Set", "Finset",
        "Bool", "Prop", "Type", "Sort",
    })
    _CAP_IDENT_RE = re.compile(r"[A-Z][A-Za-z0-9_']*")

    def names_a_hallucinated_symbol(tactic):
        if not isinstance(tactic, str):
            return []
        found = []
        for match in _CAP_IDENT_RE.findall(tactic):
            if match not in STDLIB_NAMES and match not in found:
                found.append(match)
        return found
\end{verbatim}

\noindent
\textbf{Catches:} an invented dotted path whose root is capitalised, such as
\texttt{Submodule.eq\_top\_of\_orthogonal\_eq\_bot}, while whitelisting the ten
library roots so that \texttt{exact Nat.add\_comm \_ \_} is not misread as an
invention.

\noindent
\textbf{Misses:} a great deal, and asymmetrically. Lean says ``unknown
identifier'' only for a name that resolves to nothing; a model that invents a
\emph{lowercase} lemma name --- which is the Mathlib convention for theorems,
and therefore the common case --- is caught only if the tactic happens to
mention some other capitalised token. A hallucinated name that collides with a
real declaration produces a type error instead and is never seen. The reported
rate is a floor, not an estimate.

Note also that the whitelist is a list of \emph{roots}, but the regex matches
every capitalised token in the tactic. In
\texttt{simpa using (Group.orderOf\_conj ...)} the trigger is \texttt{Group},
which is a perfectly real Mathlib class; the invention is the lowercase
\texttt{orderOf\_conj} that follows it. The rule fires for the right proof for
the wrong reason.

\subsection{SYNTAX\_ERROR\_LOOP}

\begin{verbatim}
    syntax = [s for s in steps
              if "expected" in _joined(s) or "parse error" in _joined(s)]
    if len(syntax) >= 3:
        return "SYNTAX_ERROR_LOOP", syntax[:3]
\end{verbatim}

\noindent
\textbf{Catches:} an agent emitting text Lean cannot parse at all --- prose,
truncated terms, JSON fragments that escaped the tactic extractor.

\noindent
\textbf{Misses:} the threshold is three \emph{anywhere} in the proof, not three
consecutively, so this is a rate test wearing the name of a loop. It also
competes with \texttt{ERROR\_LOOP}, which is tested first: an agent that emits
the same unparseable string three times in a row is classified as an error loop,
not a syntax loop. Only agents that fail to parse in three \emph{different} ways
reach this rule.

\subsection{BUDGET\_EXHAUSTION\_MIXED}

\begin{verbatim}
    return "BUDGET_EXHAUSTION_MIXED", steps[-3:]
\end{verbatim}

\noindent
\textbf{Catches:} everything else.

\noindent
\textbf{Misses:} nothing, which is the problem. \S\ref{sec:ft_negative}.

\section{Provider Failure Filtering}
\label{sec:ft_provider}

\subsection{The problem}

The detection rules read Lean messages, but the \texttt{lean\_messages} field
carries whatever the harness recorded for that step --- including the reason a
step never reached Lean. When NVIDIA NIM returned 404 for twenty minutes, the
affected steps recorded

\begin{verbatim}
    LLM call failed: NotFoundError: Error code: 404
\end{verbatim}

\noindent identically, step after step. That satisfies the \texttt{ERROR\_LOOP}
test perfectly. Classified naively, an outage is indistinguishable from a model
that will not stop repeating itself, and the taxonomy reports an infrastructure
event as a reasoning defect.

\subsection{What the filter actually is}

The filter is \emph{not} a keyword search over \texttt{lean\_messages}.
It is an exclusion on the proof's terminal result code, which the harness
assigns when the attempt ends:

\begin{verbatim}
    PROVIDER_CODES = {"FAILED_LLM_ERROR",
                      "FAILED_CONNECTIVITY_OUTAGE",
                      "FAILED_ELABORATION"}

    if proof.result in PROVIDER_CODES:
        continue            # not the model's failure; do not classify
\end{verbatim}

This distinction matters, because a keyword heuristic over the messages would
not have worked. Counting first-messages across both Easy-50 runs, the string
\texttt{RateLimitError: Error code: 429} occurs 165 times inside proofs that
ended in a provider result code --- and 193 times inside proofs that ended in
\texttt{FAILED\_BUDGET\_EXHAUSTED}. The same marker is more common in the set it
would have to keep than in the set it would have to drop. Rate limiting is
pervasive rather than terminal: the harness retries, the retry succeeds, the
proof continues and eventually runs out of steps for reasons of its own.
Filtering on result codes works because the harness has already decided, at the
end of the attempt, what killed it.

\subsection{Before and after}

Table~\ref{tab:ft_filter} gives the distribution under three successively
stricter definitions of ``the model failed''.
\emph{Raw} classifies every failed proof.
\emph{Result-coded} drops proofs whose terminal code blames the provider; this
is the column Chapter~\ref{chap:results} reports.
\emph{Strict} additionally drops proofs in which provider errors consumed a
majority of the steps --- rate-limit casualties that the circuit breaker did not
catch before the step budget ran out.

\begin{table}[htbp]
\centering
\caption{Failure mode distribution under three filters. Excluding provider
failures cuts \texttt{ERROR\_LOOP} from 28 to 6; the remaining shape is stable.}
\label{tab:ft_filter}
\begin{tabular}{@{}lrrrrrr@{}}
\hline
 & \multicolumn{2}{c}{\textbf{Raw}} & \multicolumn{2}{c}{\textbf{Result-coded}}
 & \multicolumn{2}{c}{\textbf{Strict}} \\
\textbf{Mode} & n & \% & n & \% & n & \% \\
\hline
\texttt{NO\_PROGRESS\_LOOP}        &  0 &  0.0 &  0 &  0.0 &  0 &  0.0 \\
\texttt{ERROR\_LOOP}               & 28 & 32.9 &  6 & 17.1 &  6 & 28.6 \\
\texttt{HALLUCINATED\_LEMMA}       &  2 &  2.4 &  2 &  5.7 &  1 &  4.8 \\
\texttt{SYNTAX\_ERROR\_LOOP}       &  2 &  2.4 &  2 &  5.7 &  2 &  9.5 \\
\texttt{BUDGET\_EXHAUSTION\_MIXED} & 53 & 62.4 & 25 & 71.4 & 12 & 57.1 \\
\hline
\textbf{Total}                     & 85 & 100.0 & 35 & 100.0 & 21 & 100.0 \\
\hline
\end{tabular}
\end{table}

Every one of the 22 \texttt{ERROR\_LOOP} cases removed by the result-code filter
carried the code \texttt{FAILED\_LLM\_ERROR}. None were borderline.

The \emph{strict} column exists because the result-code filter is not sufficient.
Provider errors that occur \emph{inside} a proof which then exhausts its budget
survive it, and they are not rare: across the 35 result-coded failures, 206 of
525 steps (39.2\%) were provider errors rather than Lean verdicts. That figure is
almost entirely one run. The Nemotron failures are clean --- 9 provider-error
steps out of 225 (4.0\%), with 9 of 15 proofs entirely free of them. The Groq
failures are not: 197 of 300 steps (65.7\%), with 14 of 20 proofs spending more
than half their budget on rate limits. Fourteen of the twenty-five
\texttt{BUDGET\_EXHAUSTION\_MIXED} cases are Groq proofs that never got enough
uninterrupted turns to fail on the mathematics.

The shape survives all three filters, which is the useful conclusion: the zero
stays zero, the residual bucket stays dominant, and no reordering of the modes
occurs.

\section{BUDGET\_EXHAUSTION\_MIXED as a Negative Result}
\label{sec:ft_negative}

Between 57\% and 71\% of genuine failures, depending on the filter, land in the
bucket defined as ``none of the above''. A taxonomy whose residual category is
its largest is not describing its domain.

It is worth being precise about what those proofs look like, because the bucket
is not homogeneous and its name is misleading. The agent in these proofs is
typically not stuck in any detectable sense. It proposes a tactic, Lean rejects
it, it proposes a different tactic, Lean rejects that too, and after fifteen
such rounds the budget ends. Sometimes real progress happens first --- 26 of the
35 result-coded failures contain at least one \texttt{advanced} step --- and
then the agent stalls at a subgoal it cannot close. Nothing repeats often enough
to trip a loop detector, nothing is obviously invented, nothing fails to parse
three times. The proof simply does not get finished.

That pattern deserves a name, and the taxonomy does not have one. A sixth mode
--- \textsc{Partial Progress Stall}, for a proof that recorded at least one
\texttt{advanced} step, never repeated a tactic more than twice, and terminated
on budget --- would separate the genuinely instructive failures from the
``never got started'' cases currently pooled with them. On this data it would
claim a substantial share of the residual and leave a much smaller true
remainder.

Three implications follow, in increasing order of interest.

\paragraph{The step budget is miscalibrated.}
Chapter~\ref{chap:results} found no proof closing after step~9 and 33 of 38
closing within four. A budget of 15 therefore spends most of its steps in a
region where success has never once been observed. Cutting it to 6 would cost
two proofs out of 38 and return roughly 60\% of the compute currently burned on
failures. For a study rate-limited by its providers, that is the single
cheapest available improvement.

\paragraph{Early stopping should key on progress, not on step count.}
A proof that has made no \texttt{advanced} step by step~5 is in a different
state from one that has decomposed its goal twice and is working on a subgoal.
The first should be abandoned; the second is exactly where extra budget might
pay. The current harness cannot tell them apart because it stops on step count
alone.

\paragraph{The residual is where the interesting failures are.}
The four specific modes describe an agent that is malfunctioning --- looping,
inventing, emitting garbage. The residual describes an agent that is behaving
sensibly and still failing, which is the harder and more relevant condition.
The taxonomy was built to catch pathology and the dominant behaviour is not
pathological; it is competent, patient, and insufficient.

\section{NO\_PROGRESS\_LOOP = 0}
\label{sec:ft_zero}

\subsection{The hypothesis}

The study was designed expecting \texttt{NO\_PROGRESS\_LOOP} to be the dominant
mode, on the order of 40\% of failures. The reasoning was ordinary: a greedy
agent with no backtracking, shown the same goal repeatedly, would propose the
same tactic repeatedly and spin until its budget ran out.

The observed count is zero. Not small --- zero, in both Easy-50 runs, under both
models, across all 20 run directories on disk, and under all three filters in
Table~\ref{tab:ft_filter}.

\subsection{Why}

The mode presupposes a \emph{silent} failure: a tactic Lean accepts while
leaving the goal unchanged. That is the only way the \texttt{no\_progress}
outcome is produced. In practice the agent's bad tactics are not accepted
silently --- they are rejected loudly. Lean returns a message, the harness
appends it to the history, and the next prompt contains it. The agent is never
in the epistemic position the hypothesis assumed.

This is best read as a property of the environment rather than of the agent, and
a favourable one. The Lean REPL does not silently tolerate a wrong tactic. Every
step returns a verdict: the goal closed, the goal changed, or an error
explaining what went wrong. An agent working against such an environment cannot
be confused about whether its last action accomplished anything. Whatever else
is hard about tactic-by-tactic proving, the credit assignment is free.

\subsection{What it invalidates}

An entire family of proposed improvements addresses this non-condition:
goal-state diffing to detect stagnation, explicit no-progress counters, forced
tactic diversity after a repeat, prompt instructions to avoid repeating
ineffective tactics. On this evidence all of it targets a failure that does not
occur. The agent's difficulty is not that it fails to notice failure --- it is
that noticing failure tells it nothing about what to do instead. The 36.4\%
tactic repetition rate reported in Chapter~\ref{chap:results} is repetition
\emph{after an explicit error message}, which is a different and harder defect
than blind spinning.

\section{Representative Traces}
\label{sec:ft_traces}

Traces are quoted from the logs with whitespace normalised and long Lean
messages truncated. The reliability caveat in \S\ref{sec:ft_threats} applies to
the first example and is stated there.

\subsection{HALLUCINATED\_LEMMA}

Theorem \texttt{TOP\_E\_01}: given that preimages of open sets are open, show
\texttt{Continuous f}. Nemotron, run \texttt{20260909\_002719},
\texttt{proof\_0008}:

\begin{verbatim}
  step 1   exact continuous_iff_isOpen.mpr h
           Unknown identifier `continuous_iff_isOpen.mpr`
  step 2   apply continuous_iff_isOpen.mpr
  step 5   exact continuous_iff_isOpen.mpr h
  step 8   exact continuous_iff_isOpen.mpr h
  step 12  exact continuous_iff_isOpen.mpr h
  (15 steps, FAILED_BUDGET_EXHAUSTED)
\end{verbatim}

The detector reports the invented symbol as \texttt{Open}, the capitalised
token inside \texttt{continuous\_iff\_isOpen}; the actual invention is the whole
name. Mathlib has no \texttt{continuous\_iff\_isOpen}. The lemma the model
wanted exists and is called \texttt{continuous\_def}:

\begin{verbatim}
  theorem continuous_def {_ : TopologicalSpace X} {_ : TopologicalSpace Y}
      {f : X -> Y} : Continuous f <-> forall s, IsOpen s -> IsOpen (f ^-1' s)
\end{verbatim}

This is the characteristic shape of the mode. The model has the mathematics
exactly right --- it knows the goal follows from an iff-characterisation of
continuity in terms of open preimages, and \texttt{.mpr} is the correct
direction to apply. What it does not have is the name. It generates a
plausible Mathlib-style identifier from the semantics, which is precisely what a
language model trained on naming conventions would do, and the identifier does
not exist.

\subsection{ERROR\_LOOP}

Theorem \texttt{SET\_E\_03}, transitivity of subset. Groq, run
\texttt{20260913\_165146}, \texttt{proof\_0006}, with no provider errors at all:

\begin{verbatim}
  step 10  exact h1.trans h2        Lean error: No goals to be solved
  step 11  exact h1.trans h2        Lean error: No goals to be solved
  step 12  done                     (accepted)
  step 15  exact h1.trans h2        Lean error: No goals to be solved
  (15 steps, FAILED_BUDGET_EXHAUSTED)
\end{verbatim}

The rule fires on steps 10, 11 and 15 --- the last three message-producing
steps, all carrying byte-identical text, with the silent step 12 skipped by the
filter. The model repeats one tactic verbatim across a span of six steps.

The honest reading of this trace is not flattering to the taxonomy. Four of the
six result-coded \texttt{ERROR\_LOOP} cases repeat the message
\texttt{No goals to be solved}, and a fifth repeats \texttt{Unknown proof
state.} after a REPL timeout. Both are messages about the \emph{proof state},
not about the tactic. Only one of the six --- \texttt{proof\_0012}, repeating
\texttt{unknown tactic} while the model emitted English prose --- is an error
loop in the sense the mode was written to capture. \texttt{ERROR\_LOOP} as
implemented is mostly detecting sessions whose state has become incoherent.

\subsection{BUDGET\_EXHAUSTION\_MIXED}

Theorem \texttt{exercise\_4\_4}, on roots of a complex polynomial and its
derivative. Groq, run \texttt{20260913\_165146}, \texttt{proof\_0018}:

\begin{verbatim}
  steps 1-4, 6-10   LLM call failed: RateLimitError: 429   (9 steps lost)
  step 5   exact Iff.rfl         Type mismatch: Iff.rfl has type ?m <-> ?m
  step 11  apply Iff.intro       (accepted, goal split)
  step 12  intro hdeg            (accepted, hypothesis introduced)
  step 13  simpa [hdeg]          Tactic `assumption` failed
  step 14  simp [hdeg]           (accepted, goal simplified)
  step 15  LLM call failed: RateLimitError: 429
  final goal:  case mp ... hdeg : p.degree = sorry |- sorry () = bot
\end{verbatim}

This is the dominant mode in one trace. Steps 11, 12 and 14 are real progress:
the iff is split, the hypothesis is introduced, the goal is simplified. Nothing
repeats, nothing is invented, nothing fails to parse more than once. The agent
was working correctly and ran out of room --- ten of its fifteen steps went to
rate limiting, and the four it got after step 10 were spent productively.

The final goal also shows the dataset problem from another angle. The statement
contains \texttt{sorry} placeholders where \texttt{card} failed to elaborate, so
the agent was working on a goal that is not the theorem anyone intended. It made
genuine progress on a malformed target.

\section{Threats to the Classification}
\label{sec:ft_threats}

\paragraph{One run predates a harness fix.}
Until commit \texttt{dd33000} (2026-09-12) the harness decided a step had failed
by searching Lean's messages for the literal word ``error''. Lean reports
\texttt{Unknown identifier} and \texttt{Unknown constant} without using it, so
failed tactics were recorded with outcome \texttt{advanced} and the harness
adopted the broken proof state they returned --- which is why so many
subsequent steps report \texttt{No goals to be solved}. The Nemotron run
(2026-09-09) predates the fix; the Groq run (2026-09-13) does not.
Re-scoring the Nemotron logs with the corrected rule reclassifies 13 of 316
steps and drops the step-1 \texttt{advanced} rate from 23.9\% to 6.5\%.
Mode assignments are less affected than step outcomes, because the rules key on
message text rather than on the outcome label --- with the exception of
\texttt{NO\_PROGRESS\_LOOP}, whose count is zero under either labelling.
Traces in \S\ref{sec:ft_traces} were drawn from the post-fix run wherever
possible; the \texttt{HALLUCINATED\_LEMMA} example is pre-fix, but the Lean
message it turns on is unaffected by the bug.

\paragraph{At least one proof was scored as a failure after closing its goal.}
In \texttt{proof\_0006} above, step 3 applied
\texttt{exact fun x hx => h2 (h1 hx)} --- a correct and complete proof of
\texttt{SET\_E\_03} --- after which no goals remained, as the subsequent
\texttt{No goals to be solved} messages confirm. The harness recorded the step
as \texttt{advanced} rather than \texttt{proved} and the attempt continued to
budget exhaustion. Six other proofs show the same shape, but in those the
closing tactic was \texttt{sorry}, \texttt{admit} or \texttt{done}, which
correctly do not count. This one appears to be a genuine proof scored as a
failure, which would make the Groq headline 9/50 rather than 8/50. It has not
been corrected in Chapter~\ref{chap:results} pending confirmation.

\paragraph{The counts are small.}
Twenty-one to thirty-five classified failures, spread over five modes, three of
which hold fewer than seven proofs each. No claim about the relative frequency
of the minority modes is supported by this data. The two findings that are
robust are the ones that do not depend on small counts: the zero, and the
dominance of the residual.
